% GENERATED FILE. Edit canonical lessons, then run npm run generate:books.
% canonical-source-hash: fnv1a64:e2546f0be49fbe41
% canonical-lessons: JA-C175-ryouri, JA-C175-yuderu, JA-C175-mazeru, JA-C175-chawan, JA-C175-abura

\chapter{Cooking at Home}
\label{ch:ja-cooking-at-home}
\hlchaptermodality{175}

\begin{chapteropening}
\textbf{By the end of this chapter:} \emph{I can say how I cook at home.}

Complete the last lesson of chapter 175: I can say how I cook at home.
\end{chapteropening}

\section[\texorpdfstring{\ja{りょうり} (ryōri)}{りょうり (ryōri)}]{\ja{りょうり} (ryōri) — cooking; a dish, cuisine}

\label{lesson:JA-C175-ryouri}



\begin{quote}
Before the new one: say the Japanese for plate, then the Japanese for to dish up.
\end{quote}

\subsection*{You'll want to know: \ja{りょうり}}
\textbf{\ja{りょうり}} — \emph{ryōri} — \textquotedblleft{}cooking; a dish, cuisine\textquotedblright{}.

\begin{quote}
\textbf{\ja{あには} \ja{りょうりが} \ja{うまいです。}} — \emph{ani wa ryōri ga umai desu} — \textquotedblleft{}My older brother is good at cooking.\textquotedblright{}
\end{quote}

\begin{cousinweb}
\textbf{\ja{りょうり}} (\emph{ryōri}) is Sino-Japanese (kango): two Chinese characters read with their on'yomi, readings borrowed from Chinese. They are \emph{ryō}, \textquotedblleft{}measure, materials\textquotedblright{} (Middle Chinese \emph{lew}) and \emph{ri}, \textquotedblleft{}arrange, reason\textquotedblright{} (\emph{liX}) (KANJIDIC2; Middle Chinese in Baxter's transcription).
\end{cousinweb}

\begin{grammarlens}[title={What you've built}]
One more word for everyday life at A2.
\end{grammarlens}

\subsection*{Your turn}
\begin{itemize}
\raggedright
  \item \emph{Say it:} \emph{ryōri}
  \item \emph{Say it:} \emph{ryōri}, once more
  \item \emph{Say it:} \emph{ani wa ryōri ga umai desu} — My older brother is good at cooking.
  \item \emph{Recall:} say the Japanese for plate, then the Japanese for to dish up, then say \emph{ryōri} again
\end{itemize}

\begin{tcolorbox}[breakable,colback=teal!4,colframe=teal!35!black,title={Before you move on}]
What does \ja{りょうり} mean? (\textquotedblleft{}Cooking; a dish, cuisine\textquotedblright{}.) Say it once more.
\end{tcolorbox}
\section[\texorpdfstring{\ja{ゆでる} (yuderu)}{ゆでる (yuderu)}]{\ja{ゆでる} (yuderu) — to boil food; to cook in hot water}

\label{lesson:JA-C175-yuderu}



\begin{quote}
Before the new one: say the Japanese for to dish up, then the Japanese for cooking.
\end{quote}

\subsection*{You'll want to know: \ja{ゆでる}}
\textbf{\ja{ゆでる}} — \emph{yuderu} — \textquotedblleft{}to boil food; to cook in hot water\textquotedblright{}.

\begin{quote}
\textbf{\ja{あさ} \ja{たまごを} \ja{ゆでます。}} — \emph{asa tamago o yudemasu} — \textquotedblleft{}I boil eggs in the morning.\textquotedblright{}
\end{quote}

Polite \textbf{\ja{ゆでます}}, \emph{yudemasu}; an ichidan verb.

\begin{cousinweb}
\textbf{\ja{ゆでる}} (\emph{yuderu}) is a native Japanese word (wago), inherited rather than borrowed; JMdict also lists a spelling with the kanji for \emph{yu}, \textquotedblleft{}hot water\textquotedblright{}, which suggests a link to that word, though the sources here do not settle it (JMdict, KANJIDIC2).
\end{cousinweb}

\begin{grammarlens}[title={What you've built}]
One more word for everyday life at A2.
\end{grammarlens}

\subsection*{Your turn}
\begin{itemize}
\raggedright
  \item \emph{Say it:} \emph{yuderu}
  \item \emph{Say it:} \emph{yuderu}, once more
  \item \emph{Say it:} \emph{asa tamago o yudemasu} — I boil eggs in the morning.
  \item \emph{Recall:} say the Japanese for to dish up, then the Japanese for cooking, then say \emph{yuderu} again
\end{itemize}

\begin{tcolorbox}[breakable,colback=teal!4,colframe=teal!35!black,title={Before you move on}]
What does \ja{ゆでる} mean? (\textquotedblleft{}To boil food; to cook in hot water\textquotedblright{}.) Say it once more.
\end{tcolorbox}
\section[\texorpdfstring{\ja{まぜる} (mazeru)}{まぜる (mazeru)}]{\ja{まぜる} (mazeru) — to mix; to stir, to blend}

\label{lesson:JA-C175-mazeru}



\begin{quote}
Before the new one: say the Japanese for cooking, then the Japanese for to boil food.
\end{quote}

\subsection*{You'll want to know: \ja{まぜる}}
\textbf{\ja{まぜる}} — \emph{mazeru} — \textquotedblleft{}to mix; to stir, to blend\textquotedblright{}.

\begin{quote}
\textbf{\ja{さとうと} \ja{しおを} \ja{まぜます。}} — \emph{satō to shio o mazemasu} — \textquotedblleft{}I mix sugar and salt.\textquotedblright{}
\end{quote}

Polite \textbf{\ja{まぜます}}, \emph{mazemasu}; an ichidan verb.

\begin{cousinweb}
\textbf{\ja{まぜる}} (\emph{mazeru}) is a native Japanese word (wago), inherited rather than borrowed from Chinese; its kanji was attached later as a kun'yomi (JMdict, KANJIDIC2).
\end{cousinweb}

\begin{grammarlens}[title={What you've built}]
One more word for everyday life at A2.
\end{grammarlens}

\subsection*{Your turn}
\begin{itemize}
\raggedright
  \item \emph{Say it:} \emph{mazeru}
  \item \emph{Say it:} \emph{mazeru}, once more
  \item \emph{Say it:} \emph{satō to shio o mazemasu} — I mix sugar and salt.
  \item \emph{Recall:} say the Japanese for cooking, then the Japanese for to boil food, then say \emph{mazeru} again
\end{itemize}

\begin{tcolorbox}[breakable,colback=teal!4,colframe=teal!35!black,title={Before you move on}]
What does \ja{まぜる} mean? (\textquotedblleft{}To mix; to stir, to blend\textquotedblright{}.) Say it once more.
\end{tcolorbox}
\section[\texorpdfstring{\ja{ちゃわん} (chawan)}{ちゃわん (chawan)}]{\ja{ちゃわん} (chawan) — rice bowl; teacup}

\label{lesson:JA-C175-chawan}



\begin{quote}
Before the new one: say the Japanese for to boil food, then the Japanese for to mix.
\end{quote}

\subsection*{You'll want to know: \ja{ちゃわん}}
\textbf{\ja{ちゃわん}} — \emph{chawan} — \textquotedblleft{}rice bowl; teacup\textquotedblright{}.

\begin{quote}
\textbf{\ja{ちゃわんに} \ja{ごはんを} \ja{もります。}} — \emph{chawan ni gohan o morimasu} — \textquotedblleft{}I dish up rice into a bowl.\textquotedblright{}
\end{quote}

\begin{cousinweb}
\textbf{\ja{ちゃわん}} (\emph{chawan}) is Sino-Japanese (kango): two Chinese characters read with their on'yomi, readings borrowed from Chinese. They are \emph{cha}, \textquotedblleft{}tea\textquotedblright{} (Middle Chinese \emph{drae}) and \emph{wan}, \textquotedblleft{}bowl\textquotedblright{} (\emph{'wanX}) (KANJIDIC2; Middle Chinese in Baxter's transcription).
\end{cousinweb}

\begin{grammarlens}[title={What you've built}]
One more word for everyday life at A2.
\end{grammarlens}

\subsection*{Your turn}
\begin{itemize}
\raggedright
  \item \emph{Say it:} \emph{chawan}
  \item \emph{Say it:} \emph{chawan}, once more
  \item \emph{Say it:} \emph{chawan ni gohan o morimasu} — I dish up rice into a bowl.
  \item \emph{Recall:} say the Japanese for to boil food, then the Japanese for to mix, then say \emph{chawan} again
\end{itemize}

\begin{tcolorbox}[breakable,colback=teal!4,colframe=teal!35!black,title={Before you move on}]
What does \ja{ちゃわん} mean? (\textquotedblleft{}Rice bowl; teacup\textquotedblright{}.) Say it once more.
\end{tcolorbox}
\section[\texorpdfstring{\ja{あぶら} (abura)}{あぶら (abura)}]{\ja{あぶら} (abura) — oil; fat}

\label{lesson:JA-C175-abura}



\begin{quote}
Before the new one: say the Japanese for to mix, then the Japanese for rice bowl.
\end{quote}

\subsection*{You'll want to know: \ja{あぶら}}
\textbf{\ja{あぶら}} — \emph{abura} — \textquotedblleft{}oil; fat\textquotedblright{}.

\begin{quote}
\textbf{\ja{あぶらで} \ja{にくを} \ja{やきます。}} — \emph{abura de niku o yakimasu} — \textquotedblleft{}I cook the meat in oil.\textquotedblright{}
\end{quote}

\begin{cousinweb}
\textbf{\ja{あぶら}} (\emph{abura}) is a native Japanese word (wago), inherited rather than borrowed from Chinese; its kanji was attached later as a kun'yomi (JMdict, KANJIDIC2).
\end{cousinweb}

\begin{grammarlens}[title={What you've built}]
One more word for everyday life at A2.
\end{grammarlens}

\subsection*{Your turn}
\begin{itemize}
\raggedright
  \item \emph{Say it:} \emph{abura}
  \item \emph{Say it:} \emph{abura}, once more
  \item \emph{Say it:} \emph{abura de niku o yakimasu} — I cook the meat in oil.
  \item \emph{Recall:} say the Japanese for to mix, then the Japanese for rice bowl, then say \emph{abura} again
\end{itemize}

\begin{tcolorbox}[breakable,colback=teal!4,colframe=teal!35!black,title={Before you move on}]
What does \ja{あぶら} mean? (\textquotedblleft{}Oil; fat\textquotedblright{}.) Say it once more.
\end{tcolorbox}
